%% start of file `template.tex'.
%% Copyright 2006-2015 Xavier Danaux (xdanaux@gmail.com), 2020-2024 moderncv maintainers (github.com/moderncv).
%
% This work may be distributed and/or modified under the
% conditions of the LaTeX Project Public License version 1.3c,
% available at http://www.latex-project.org/lppl/.


\documentclass[11pt,a4paper,sans]{moderncv}        % possible options include font size ('10pt', '11pt' and '12pt'), paper size ('a4paper', 'letterpaper', 'a5paper', 'legalpaper', 'executivepaper' and 'landscape') and font family ('sans' and 'roman')

\moderncvcolor{black}
\moderncvstyle[]{banking}
\nopagenumbers{}                                  % uncomment to suppress automatic page numbering for CVs longer than one page

\usepackage{hyperref}

% adjust the page margins
\usepackage[hmargin=0.5in,vmargin=10pt]{geometry}

% font loading
% for luatex and xetex, do not use inputenc and fontenc
% see https://tex.stackexchange.com/a/496643
\ifxetexorluatex
  \usepackage{fontspec}
  \usepackage{unicode-math}
  \usepackage{FiraSans}
  \defaultfontfeatures{Ligatures=TeX}
  \setmainfont{Fira Sans}
  \setsansfont[ItalicFont={Fira Sans Light Italic},
            BoldFont={Fira Sans Bold},
            BoldItalicFont={Fira Sans Italic}]
            {Fira Sans Light}
  % \setmonofont{Latin Modern Mono}
  \setmonofont{Monaspace Neon}
  \setmathfont{Latin Modern Math}

  % you may also consider Fira Sans Light for a extra modern look
  %\setsansfont[ItalicFont={Fira Sans Light Italic},%
  %           BoldFont={Fira Sans},%
  %           BoldItalicFont={Fira Sans Italic}]%
  %           {Fira Sans Light}%
\else
  \usepackage[utf8]{inputenc}
  \usepackage[T1]{fontenc}
  \usepackage{lmodern}
\fi

% document language
\usepackage[english]{babel}  % FIXME: using spanish breaks moderncv

\usepackage{setspace}
\linespread{1.5}

% personal data
\name{\textmd{\textsc{ShengQiao}}}{\textmd{\textsc{Wu}}}
\phone[mobile]{+86~159~8583~0350}
\email{Blind-Guess-Senior@outlook.com}
\homepage{https://blind-guess-senior.github.io/}

\social[github]{Blind-Guess-Senior}

% \extrainfo{additional information}                 % optional, remove / comment the line if not wanted
% \quote{Some quote}                                 % optional, remove / comment the line if not wanted

%----------------------------------------------------------------------------------
%            content
%----------------------------------------------------------------------------------
\begin{document}
%\begin{CJK*}{UTF8}{gbsn}                          % to typeset your resume in Chinese using CJK
%-----       resume       ---------------------------------------------------------
\makecvtitle

\section{Education}  % for 'contemporary' style use optional argument for displaying an icon, e.g. \section[\faGraduationCap]{Education}
\cventry{2024 -- present}{Undergraduate | \textit{Computer Science}}{ShanghaiTech University}{\textit{Shanghai, China}}{}{}  % arguments 3 to 6 can be left empty
% \cventry{year--year}{Degree}{Institution}{City}{\textit{Grade}}{Description}

\section{Skills}
\cvitem{Languages}{C\# | C/C++ | Javascript/Typescript | OCaml | Python | LaTeX}

\section{Competitions}
\cventry{Oct 2022}{CCF Certified Software Professional Senior (CSP-S), Round 2}{CCF CSP-S, Round 2 \textmd{\textsc{| Second Prize}}}{C++, Argorithm}{}{}
\cventry{Oct 2022}{CCF Certified Software Professional Senior (CSP-S), Round 1}{CCF CSP-S, Round 1 \textmd{\textsc{| Second Prize}}}{C++, Argorithm}{}{}

\section{Games}
\cventry{Jul 2026}{\textcolor[RGB]{128,128,128}{A 2D rhythm game with a quirky input method: smashing the keyboard.}}{KBre \textmd{\textsc{| Team Project}}}{Unity}{}{
  \vspace{0.2em}
  Role: Lead Programmer; Team of 5
  \vspace{0.1em}
  \begin{itemize}
    \item Players control movement and attacks using keyboard-smashing input patterns such as simple smashes, slides, and double smashes; timing these inputs to the beat increases damage and reduces resource costs.
    \item I designed the overall code architecture and implemented the event system, input-pattern recognition system, scene object logic, scene management, and enemy state machines.
  \end{itemize}}
\vspace{0.4em}

\cventry{Jul 2025}{\textcolor[RGB]{128,128,128}{A 2D linear platformer with ability-gated progression.}}{C-17ES-α \textmd{\textsc{| Course Project}}}{Unity}{}{
  \vspace{0.2em}
  Role: Lead Programmer; Team of 5
  \vspace{0.1em}
  \begin{itemize}
    \item The player starts with only the ability to move right, gradually unlocking left, up and down movement as they progress. Each direction provides a unique special effect, and the player must strategically use these effects to overcome obstacles and complete levels.
    \item I designed the overall code architecture and implemented the event system, movement ability system, and interactive scene objects.
  \end{itemize}}
\vspace{0.4em}

\section{Projects}
\cventry{Mar 2026 -- present}{\textcolor[RGB]{128,128,128}{A compile-based, simple yet powerful dialogue library for Unity.}}{Artifact Dialoguer \textmd{\textsc{| Personal}}}{C\#, .NET, Unity}{}{
  \vspace{0.2em}
  \begin{itemize}
    \item I designed a concise and unambiguous dialogue script syntax, and implemented an LL(1) parser to compile it into Unity ScriptableObjects.
    \item I implemented a runtime controller that reads the compiled data and controls dialogue progression during gameplay.
  \end{itemize}
}

\cventry{Feb 2026 -- present}{\textcolor[RGB]{128,128,128}{A less-dependent Unity tool library.}}{Artifact Unity Utils \textmd{\textsc{| Personal}}}{C\#, .NET, Unity}{}{
  \vspace{0.2em}
  \begin{itemize}
    \item I implemented some useful design patterns and data structures in .NET, including an event-bus based event system, dependency injection, quadtree and so on.
  \end{itemize}
}
\vspace{0.4em}

\cventry{Sep 2025 -- Jan 2026}{\textcolor[RGB]{128,128,128}{A simple language compiler for the OATC language.}}{OATC Language Compiler \textmd{\textsc{| Personal}}}{OCaml, LLVM, X86lite}{}{
  \vspace{0.2em}
  \begin{itemize}
    \item I implemented an X86lite instruction set simulator and assembler. And the OATC language interpreter using OCaml.
    \item I've also developed the compiler from LLVMlite IR to X86lite platform with frontend from OATC to LLVMlite IR.
  \end{itemize}
}
\vspace{0.4em}

\cventry{Feb 2025 -- Jun 2025}{\textcolor[RGB]{128,128,128}{An operating system for the 80x86 architecture.}}{PintOS \textmd{\textsc{| Personal}}}{C, x86 Assembly}{}{
  \vspace{0.2em}
  \begin{itemize}
    \item I implemented the advanced scheduling, system call, user/kernel mode, virtual memory and file system of the operating system.
    \item PintOS contains basic shell and filesystem, and is able to run programs in user mode.
  \end{itemize}
}
\vspace{0.4em}
% \vspace{0.5em}

% Detailed achievements:
% \begin{itemize}
% \item Achievement 1
% \item Achievement 2 (with sub-achievements)
%   \begin{itemize}
%   \item Sub-achievement (a);
%   \item Sub-achievement (b), with sub-sub-achievements (don't do this!);
%     \begin{itemize}
%     \item Sub-sub-achievement i;
%     \item Sub-sub-achievement ii;
%     \item Sub-sub-achievement iii;
%     \end{itemize}
%   \item Sub-achievement (c);
%   \end{itemize}
% \item Achievement 3
% \item Achievement 4
% \end{itemize}}
% \cventry{year--year}{Job title}{Employer}{City}{}{Description line 1\newline{}Description line 2\newline{}Description line 3}
% \subsection{Miscellaneous}
% \cventry{year--year}{Job title}{Employer}{City}{}{Description}

\clearpage

\end{document}


%% end of file `template.tex'.
